\chapter{闭区间上的连续函数}
\label{chap:continuous-functions-closed-interval}
\label{chap:compact-connected-real-line}


连续性给出的邻域通常随点而变，单靠逐点条件不能直接产生整个定义域上的共同界。闭区间却允许从任意开覆盖中选出有限子覆盖，因而能把无穷多个局部估计缩减为有限多个。另一方面，区间不能分成两个彼此分离的非空相对开集，这一性质保证连续函数不会跳过中间值。有限覆盖性质将给出有界性和极值定理，区间的连通性则给出介值定理。


\section{开覆盖与有限子覆盖的动机}

\begin{definition}[开覆盖]
设 \(K\subset\R\)。一族开集 \(\{U_\lambda\}_{\lambda\in\Lambda}\) 若满足
\[
 K\subset\bigcup_{\lambda\in\Lambda}U_\lambda,
\]
则称它是 \(K\) 的一个\term{开覆盖}。若存在有限个指标 \(\lambda_1,\ldots,\lambda_m\) 仍使
\[
 K\subset U_{\lambda_1}\cup\cdots\cup U_{\lambda_m},
\]
则称原覆盖有\term{有限子覆盖}。
\end{definition}

连续性在每个 \(x\in K\) 处给出一个可能依赖 \(x\) 的邻域 \(U_x\)。若能从 \(\{U_x:x\in K\}\) 中选出有限多个覆盖整个 \(K\)，便可对有限组局部常数取最大值、最小值，把局部估计转成全局估计。

这里的困难不是每个点缺少局部控制，而是局部控制可能随点无限退化。有限子覆盖把无穷多个局部选择压缩成有限多个，有限集合上的最小半径和最大常数都确实存在。这正是闭区间上许多整体结论使用同一个几何事实的原因。

\begin{counterexample}
开区间 \((0,1)\) 的覆盖
\[
 U_n=(1/n,1),\qquad n\ge2,
\]
没有有限子覆盖。任取有限多个，最大指标记为 \(N\)，其并集仍不包含 \((0,1/N]\) 中的点。这提示缺失端点会破坏有限覆盖性质。
\end{counterexample}

\begin{counterexample}
实线的覆盖 \(U_n=(-n,n)\)（\(n\in\Npos\)）没有有限子覆盖，说明无界性也是障碍。
\end{counterexample}

\section{闭区间的有限覆盖定理}

下面采用“可从左端覆盖到哪里”的上确界证明。集合 \(E\) 记录已经能用有限多个覆盖成员处理的右端点；若其上确界尚未到达 \(b\)，覆盖上确界的某个开集便能把有限覆盖向右推进，从而造成矛盾。证明末尾还需单独说明上确界自身可以覆盖到，不能把 \(\sup E=b\) 与 \(b\in E\) 混为一谈。

\begin{theorem}[Heine--Borel 有限覆盖定理]
闭区间 \([a,b]\) 的每个开覆盖都有有限子覆盖。
\end{theorem}

\begin{proof}
若 \(a=b\)，任取一个包含 \(a\) 的覆盖成员即可。以下设 \(a<b\)。给定开覆盖 \(\mathcal U\)，定义
\[
 E=\{x\in[a,b]:[a,x]\text{ 可由 }\mathcal U
 \text{ 中有限多个成员覆盖}\}.
\]
因为某个 \(U\in\mathcal U\) 含 \(a\)，且 \(U\) 开，存在 \(\eta>0\) 使
\((a-\eta,a+\eta)\subset U\)，故 \(E\) 非空。由 \(E\) 的定义，每个 \(x\in E\) 都满足 \(a\le x\le b\)，因而 \(E\subset[a,b]\) 且以 \(b\) 为上界。由完备性可令 \(c=\sup E\)。

若 \(c<b\)，选择 \(V\in\mathcal U\) 使 \(c\in V\)。因 \(V\) 开，存在 \(\rho>0\) 使
\[
 (c-\rho,c+\rho)\subset V,
\]
并可缩小 \(\rho\) 使 \(c+\rho/2\le b\)。由 \(c=\sup E\)，存在 \(x\in E\) 满足 \(c-\rho/2<x\le c\)。区间 \([a,x]\) 已有有限子覆盖，再添 \(V\) 就覆盖 \([a,c+\rho/2]\)，所以 \(c+\rho/2\in E\)，与 \(c\) 是上界矛盾。故 \(c=b\)。

最后还需证明 \(b\in E\)，不能只由 \(\sup E=b\) 直接断言。取含 \(b\) 的 \(W\in\mathcal U\) 及 \(\sigma>0\)，使 \((b-\sigma,b+\sigma)\subset W\)。由上确界定义选 \(x\in E\) 且 \(b-\sigma<x\le b\)。覆盖 \([a,x]\) 的有限成员再加 \(W\) 即覆盖 \([a,b]\)。因此 \(b\in E\)，定理得证。
\end{proof}

\begin{remark}
证明三次用到结构条件：开性把覆盖成员扩成邻域；闭区间的端点属于集合；实数完备性保证 \(\sup E\) 存在。若去掉有界或闭端点，前节反例表明结论可能失败。
\end{remark}

\begin{theorem}[实线上的 Heine--Borel 刻画]
集合 \(K\subset\R\) 的每个开覆盖都有有限子覆盖，当且仅当 \(K\) 闭且有界。
\end{theorem}

\begin{proof}
若 \(K\) 闭且有界，则 \(K\subset[-M,M]\)。在覆盖中加入开集 \(\R\setminus K\)，得到 \([-M,M]\) 的开覆盖；闭区间定理给出有限子覆盖，删去可能出现的 \(\R\setminus K\) 后，余下成员覆盖 \(K\)。

反之，若 \(K\) 无界，\(\{(-n,n)\}_{n\in\Npos}\) 覆盖 \(K\) 而无有限子覆盖。若 \(K\) 不闭，存在聚点 \(a\notin K\)。开集
\[
 U_n=\{x\in\R:\abs{x-a}>1/n\},\qquad n\in\Npos,
\]
覆盖 \(K\)，但任意有限子族的并仍避开 \(a\) 的某个邻域；因 \(a\) 是 \(K\) 的聚点，该并不能覆盖 \(K\)。
\end{proof}

\section{连续函数在闭区间上有界}

\begin{theorem}[有界性定理]
若 \(f:[a,b]\to\R\) 连续，则 \(f\) 在 \([a,b]\) 上有界。
\end{theorem}

\begin{proof}[覆盖证明]
对每个 \(x\in[a,b]\)，由连续性取开邻域 \(U_x\)，使
\[
 y\in U_x\cap[a,b]\Longrightarrow\abs{f(y)-f(x)}<1.
\]
这些 \(U_x\) 构成 \([a,b]\) 的开覆盖。取有限子覆盖
\[
 U_{x_1},\ldots,U_{x_m}.
\]
令
\[
 M=1+\max_{1\le j\le m}\abs{f(x_j)}.
\]
任取 \(y\in[a,b]\)，它属于某个 \(U_{x_j}\)，故
\[
 \abs{f(y)}\le\abs{f(y)-f(x_j)}+\abs{f(x_j)}<M.
\]
\end{proof}

\begin{proof}[序列证明]
若 \(f\) 无界，可对每个 \(n\) 选 \(x_n\in[a,b]\) 使 \(\abs{f(x_n)}>n\)。由 Bolzano--Weierstrass 定理，\((x_n)\) 有收敛子列 \(x_{n_k}\to x\in[a,b]\)；闭性保证极限仍在区间中。连续性给出
\[
 f(x_{n_k})\to f(x),
\]
故像子列有界，与 \(\abs{f(x_{n_k})}>n_k\to\infty\) 矛盾。
\end{proof}

第一个证明从有限子覆盖取得有限组局部界，第二个证明从有界点列抽取收敛子列。在实数轴上，这两种论证分别对应紧致性的覆盖形式和序列形式；它们都以闭且有界为特征。

\section{最大值最小值定理}

有界性只产生数 \(M=\sup f([a,b])\)，并不保证存在 \(x\) 使 \(f(x)=M\)。取得最大值还需要两步：从上确界附近选出一列函数值，再从相应点列抽取仍留在闭区间中的收敛子列。连续性最后把点列极限送到函数值极限。

\begin{theorem}[Weierstrass 极值定理]
若 \(f:[a,b]\to\R\) 连续，则存在 \(x_{\min},x_{\max}\in[a,b]\)，使
\[
 f(x_{\min})\le f(x)\le f(x_{\max})
 \qquad(x\in[a,b]).
\]
\end{theorem}

\begin{proof}
由有界性，值集
\[
 S=f([a,b])
\]
有上界。令 \(M=\sup S\)。按上确界定义，对每个 \(n\) 选 \(x_n\in[a,b]\) 使
\[
 M-\frac1n<f(x_n)\le M.
\]
由 Bolzano--Weierstrass 定理，某子列 \(x_{n_k}\to x_{\max}\in[a,b]\)。连续性与夹逼给出
\[
 f(x_{\max})
 =\lim_{k\to\infty}f(x_{n_k})
 =M.
\]
故最大值取得。对 \(-f\) 应用同一结论，得到最小值。
\end{proof}

\begin{counterexample}[缺少条件时]
在 \((0,1)\) 上，连续函数 \(f(x)=x\) 有界但不取得最大值和最小值；在闭集 \([0,\infty)\) 上，连续函数 \(f(x)=x\) 无界；在闭区间上，函数
\[
 f(0)=0,\qquad f(x)=1/x\ (x>0)
\]
不连续且无界。闭、有界与连续分别承担不同作用。
\end{counterexample}

\section{区间的连通性}

有限覆盖性质负责把局部估计变成有限控制，连通性处理另一类问题：连续函数的像能否越过某些数值。区间性质说明两个点之间没有缺口；连续映射若把端点送到某个数值的两侧，就不能在像中删掉该中间值。

\begin{definition}[相对分离与连通]
设 \(E\subset\R\)。若存在 \(E\) 中两个非空相对开集 \(A,B\)，满足
\[
 E=A\cup B,\qquad A\cap B=\varnothing,
\]
则称 \(E\) 可分离。不能这样分离的集合称为\term{连通集}。
\end{definition}

\begin{theorem}[实线连通集的刻画]
非空集合 \(I\subset\R\) 连通，当且仅当它具有区间性质：
\[
 x,y\in I,\quad x<z<y\Longrightarrow z\in I.
\]
因此实线中的连通集恰为各种区间和单点集。
\end{theorem}

\begin{proof}
先设 \(I\) 有区间性质。若 \(I=A\cup B\) 是分离，取 \(a\in A\)、\(b\in B\)，不妨 \(a<b\)。令
\[
 c=\sup(A\cap[a,b]).
\]
由区间性质 \(c\in I\)。若 \(c\in A\)，因 \(A\) 在 \(I\) 中相对开，存在 \(c\) 右侧的 \(I\) 中点仍属于 \(A\)，与上确界矛盾；若 \(c\in B\)，相对开性给出 \(c\) 左侧一段 \(I\) 中点属于 \(B\)，与 \(c\) 可由 \(A\cap[a,b]\) 中点从左逼近矛盾。因此不可分离。

反之，若区间性质失败，存在 \(x<z<y\)，其中 \(x,y\in I\) 而 \(z\notin I\)。则
\[
 A=I\cap(-\infty,z),\qquad B=I\cap(z,\infty)
\]
是两个非空、不交的相对开集，其并为 \(I\)，故 \(I\) 不连通。
\end{proof}

\begin{remark}
连通性不等同于“点很多”或“没有孤立点”。集合 \(\Q\) 稠密且没有孤立点，但任取无理数 \(\alpha\)，可用 \((-\infty,\alpha)\cap\Q\) 与 \((\alpha,\infty)\cap\Q\) 把它分离。
\end{remark}

\section{零点定理与介值定理}

\begin{theorem}[Bolzano 零点定理]
若 \(f:[a,b]\to\R\) 连续且
\[
 f(a)f(b)<0,
\]
则存在 \(c\in(a,b)\) 使 \(f(c)=0\)。
\end{theorem}

\begin{proof}
不妨 \(f(a)<0<f(b)\)。令
\[
 E=\{x\in[a,b]:f(x)<0\},\qquad c=\sup E.
\]
有 \(a\in E\)，且 \(b\) 是上界，所以 \(c\in[a,b]\)。若 \(f(c)<0\)，连续性使 \(c\) 右侧某邻域内函数仍为负，从而存在大于 \(c\) 的 \(E\) 中点，矛盾。若 \(f(c)>0\)，连续性使 \(c\) 左侧某邻域内函数为正，但上确界定义要求有 \(E\) 中点任意靠近 \(c\) 的左侧，也矛盾。故 \(f(c)=0\)。端点符号严格相反又保证 \(c\in(a,b)\)。
\end{proof}

\begin{theorem}[介值定理]
若 \(f:I\to\R\) 在区间 \(I\) 上连续，\(x_0,x_1\in I\)，且 \(y\) 位于 \(f(x_0)\) 与 \(f(x_1)\) 之间，则存在 \(c\) 位于 \(x_0,x_1\) 之间，使 \(f(c)=y\)。
\end{theorem}

\begin{proof}
对函数 \(g(x)=f(x)-y\) 在闭区间
\[
 [\min\{x_0,x_1\},\max\{x_0,x_1\}]
\]
上应用零点定理。若 \(y\) 等于端点函数值，直接取相应端点；否则 \(g\) 的端点值严格异号。
\end{proof}

\begin{corollary}
区间上的连续函数的值域仍是区间。特别地，若连续函数只取整数值，则它在每个区间上为常值。
\end{corollary}

\section{连续映射保持紧致性与连通性}

\begin{theorem}[连续像保持有限覆盖性质]
设 \(K\subset\R\) 具有有限覆盖性质，\(f:K\to\R\) 连续。则 \(f(K)\) 也具有有限覆盖性质。
\end{theorem}

\begin{proof}
给定 \(f(K)\) 的开覆盖 \(\{V_\lambda\}\)。由连续性，每个原像
\[
 f^{-1}(V_\lambda)
\]
在 \(K\) 中相对开，并且这些原像覆盖 \(K\)。把它们写成 \(K\cap U_\lambda\)（\(U_\lambda\) 为实线开集），由 \(K\) 的有限覆盖性质选出有限多个。对应的 \(V_\lambda\) 便覆盖 \(f(K)\)。
\end{proof}

\begin{theorem}[连续像保持连通性]
若 \(E\subset\R\) 连通，\(f:E\to\R\) 连续，则 \(f(E)\) 连通。
\end{theorem}

\begin{proof}
若 \(f(E)=A\cup B\) 可分离，则原像
\[
 f^{-1}(A),\qquad f^{-1}(B)
\]
是 \(E\) 中两个不交、非空、相对开的集合，且并为 \(E\)，这会分离 \(E\)，矛盾。
\end{proof}

有界性和极值定理也可以从连续像保持紧致性推出，介值定理则是连续像保持连通性在实数轴上的结果。以后进入一般拓扑空间时，这两个保持定理仍成立，而“闭且有界”等实数轴特有的刻画需要另行检验。

\section*{习题}
\addcontentsline{toc}{section}{习题}

\subsection*{A组　覆盖、紧致与极值}
\begin{enumerate}[label=\textbf{18.\arabic*.}]
 \exerciseitem{18.1}  分别给出 \((0,1)\)、\([0,\infty)\)、\(\R\) 没有有限子覆盖的具体开覆盖。
 \exerciseitem{18.2}  完整重建闭区间有限覆盖定理的上确界证明。
 \exerciseitem{18.3}  在上述证明中说明为何 \(\sup E=b\) 之后仍需单独证明 \(b\in E\)。
 \exerciseitem{18.4}  证明闭且有界集合 \(K\subset\R\) 具有有限覆盖性质。
 \exerciseitem{18.5}  证明具有有限覆盖性质的实数集必有界。
 \exerciseitem{18.6}  证明具有有限覆盖性质的实数集必闭。
 \exerciseitem{18.7}  分别用有限覆盖法和序列法证明连续函数在闭区间上有界。
 \exerciseitem{18.8}  完整证明 Weierstrass 极值定理。
 \exerciseitem{18.9}  为“闭”“有界”“连续”三个条件各给一个删去后极值结论失败的反例。
\end{enumerate}

\subsection*{B组　连通与介值}
\begin{enumerate}[label=\textbf{18.\arabic*.},resume]
 \exerciseitem{18.10}  证明任意区间连通。
 \exerciseitem{18.11}  证明实线中的连通集具有区间性质。
 \exerciseitem{18.12}  证明 \(\Q\) 与无理数集都不连通。
 \exerciseitem{18.13}  用上确界法证明零点定理，并注明连续性的两次使用。
 \exerciseitem{18.14}  从零点定理推出介值定理。
 \exerciseitem{18.15}  证明连续函数把区间映为区间。
 \exerciseitem{18.16}  证明区间上取值于 \(\Z\) 的连续函数必为常值。
 \exerciseitem{18.17}  若连续函数 \(f:[a,b]\to[a,b]\)，证明存在 \(c\in[a,b]\) 使 \(f(c)=c\)。
 \exerciseitem{18.18}  证明奇次实系数多项式至少有一个实根。
\end{enumerate}

\subsection*{C组　结构、边界与项目}
\begin{enumerate}[label=\textbf{18.\arabic*.},resume]
 \exerciseitem{18.19}  证明连续像保持有限覆盖性质。
 \exerciseitem{18.20}  证明连续像保持连通性。
 \exerciseitem{18.21}  证明闭区间中任意无限子集都有聚点，并由此重新证明有限覆盖定理。
 \exerciseitem{18.22}  证明紧致集上的连续双射到实数子集时逆映射连续。
 \exerciseitem{18.23} \ 【前置：\chapterref{chap:subsequences-bw}与本章；预计用时：90分钟】建立“闭且有界—序列紧致—有限覆盖性质”的三向证明链，要求每个方向独立注明使用的实数完备性形式。
 \exerciseitem{18.24}  比较有限覆盖数、Lebesgue 数与一致连续性之间的联系；可预读\chapterref{chap:uniform-continuity}后完成。
 \exerciseitem{18.25}  设 \(K_1\supset K_2\supset\cdots\) 为非空紧集。证明若 \(K_1\) 有界，则 \(\bigcap_nK_n\ne\varnothing\)。
 \exerciseitem{18.26}  证明连续函数在紧集上的零点集仍是紧集；若函数处处非零，证明其绝对值有正的最小值。
\end{enumerate}
